\chapter{Fundamente Teoretice și Studiu de Literatură}
\label{chap:literatura}

\section{Modele de învățare automată utilizate}

Predicția rezultatului unui meci de fotbal se poate formula ca o problemă
de clasificare multi-clasă: dat un \textit{snapshot}, adică o reprezentare
vectorială a stării jocului la un moment dat, conținând informații precum
scorul curent, timpul scurs, șuturile pe poartă și expected goals acumulate,
modelul trebuie să estimeze probabilitățile pentru cele trei rezultate
posibile: \textit{victorie gazdă}, \textit{egal} sau
\textit{victorie oaspeți}. În literatura de specialitate, această problemă
a fost abordată prin metode variate, de la modele liniare până la
gradient boosting și rețele Bayesiene \cite{Robberechts2019, Maddox2022}.
În cadrul acestei lucrări a fost evaluată performanța a trei algoritmi
(Gradient Boosting, LightGBM și Regresie Logistică), selectați astfel
încât să acopere un spectru de complexitate de la modelul liniar de
referință până la implementările optimizate de ansamblu. Alegerea modelului
final s-a bazat pe metrici probabilistice (RPS, Brier Score, eroare de
calibrare), nu doar pe acuratețe de clasificare, deoarece pentru aplicații
tactice calitatea distribuției de probabilitate este mai relevantă decât
simpla predicție a clasei câștigătoare \cite{Robberechts2019}.

\subsection{Gradient Boosting}

Gradient Boosting este o metodă de ansamblu care construiește secvențial o
serie de arbori de decizie, fiecare arbore nou fiind antrenat să corecteze
erorile reziduale ale celui precedent. Formal, modelul optimizează o funcție
de pierdere diferențiabilă $\mathcal{L}(y, F(x))$ prin coborâre pe gradient
în spațiul funcțiilor, adăugând la fiecare pas un estimator slab $h_t(x)$
scalat cu rata de învățare $\eta$:

\begin{equation}
    F_t(x) = F_{t-1}(x) + \eta \cdot h_t(x)
\end{equation}

unde $F_t(x)$ este modelul la iterația $t$, $F_{t-1}(x)$ este modelul
anterior, $h_t(x)$ este arborele de decizie adăugat la pasul curent, iar
$\eta \in (0, 1]$ este rata de învățare care controlează contribuția
fiecărui arbore nou.

Spre deosebire de Random Forest, care antrenează arborii în paralel și
independent, Gradient Boosting exploatează informația din pașii anteriori,
ceea ce îl face mai precis pe date structurate, dar și mai sensibil la
overfitting dacă hiperparametrii nu sunt controlați corespunzător.
Detaliile privind configurația hiperparametrilor și rezultatele obținute
sunt prezentate în Capitolul~\ref{cap:metodologie}.

\subsection{LightGBM}

LightGBM (Light Gradient Boosting Machine) este o implementare optimizată
a algoritmului Gradient Boosting, dezvoltată de Microsoft \cite{Ke2017},
care introduce două tehnici cheie față de implementările clasice:
Gradient-based One-Side Sampling (GOSS) și Exclusive Feature Bundling (EFB).

GOSS păstrează instanțele cu gradient mare (cele care contribuie cel mai
mult la procesul de învățare) și eșantionează aleatoriu instanțele cu
gradient mic, reducând volumul de date procesat fără pierdere semnificativă
de informație. EFB grupează caracteristicile mutual exclusive pentru a reduce
dimensionalitatea efectivă a datelor.

Din punct de vedere structural, LightGBM construiește arborii
\textit{leaf-wise} (extinde frunza cu cel mai mare câștig de informație),
spre deosebire de abordarea \textit{level-wise} a implementărilor
tradiționale, ceea ce duce la arbori mai adânci și mai expresivi cu același
număr de iterații. Această proprietate îl face deosebit de eficient pe
seturi de date de mari dimensiuni, cum este cel utilizat în această lucrare
\cite{StatsBomb}. Criteriile complete de selecție a acestui model și
comparația cu celelalte variante sunt detaliate în
Capitolul~\ref{cap:metodologie}.

\subsection{Regresie Logistică}

Regresia logistică este utilizată ca model de referință (baseline),
estimând probabilitățile celor trei clase prin funcția softmax:

\begin{equation}
    P(y = k \mid x) = \frac{e^{w_k^T x}}{\sum_{j=1}^{K} e^{w_j^T x}}
\end{equation}

unde $y$ este clasa prezisă, $x$ este vectorul de caracteristici, $w_k$ sunt
ponderile asociate clasei $k$, iar $K = 3$ este numărul total de clase
(\textit{victorie gazdă}, \textit{egal}, \textit{victorie oaspeți}).

Ipoteza sa fundamentală este că relația dintre caracteristici și ieșire este
liniară în spațiul log-probabilităților. Rolul său în această lucrare este
de a verifica dacă non-liniaritatea capturată de modelele de ansamblu este
cu adevărat necesară, cu alte cuvinte dacă complexitatea suplimentară
a LightGBM și Gradient Boosting se justifică prin îmbunătățiri măsurabile
ale metricilor probabilistice față de un model liniar simplu. Rezultatele
comparative sunt prezentate în Capitolul~\ref{cap:metodologie}.

Alegerea modelului final nu s-a bazat exclusiv pe acuratețea de
clasificare, ci pe metrici probabilistice precum Ranked Probability Score,
Brier Score și Expected Calibration Error, definite formal și discutate în
detaliu în Capitolul~\ref{cap:metodologie}, Secțiunea~\ref{sec:metrici},
acolo unde sunt prezentate și rezultatele experimentale care le folosesc.

% ----------------------------------------------------------------
\section{Lucrări relevante din literatură}
\label{sec:literatura}

Predicția win probability în sport are o istorie relativ scurtă ca domeniu
formalizat, dar a acumulat o literatură relevantă mai ales după
popularizarea datelor la nivel de eveniment. Prezentăm în continuare
lucrările cu cel mai direct impact asupra metodologiei adoptate în această
lucrare.

\textbf{Robberechts et al. (2019)} propun primul model de win probability
în timp real dedicat fotbalului \cite{Robberechts2019}. Autorii utilizează
date la nivel de eveniment din opt sezoane ale competițiilor europene de
top, provenite din platforma StatsBomb, acoperind aproximativ 2.000 de
meciuri. Sunt evaluate două arhitecturi: un model Gradient Boosting și o
rețea Bayesiană, ambele antrenate pe caracteristici derivate din goluri,
șuturi și presiune. Modelul Gradient Boosting obține un RPS de 0.197,
superior liniei de bază statistice, iar autorii concluzionează că
evenimentele in-game, în special golurile și șuturile pe poartă, contribuie
semnificativ la actualizarea probabilităților de-a lungul meciului.
O concluzie practică importantă este stabilirea pragului ECE de 0.05 ca
criteriu de acceptabilitate pentru uz tactic real. Această lucrare
constituie referința primară pentru metodologia adoptată în prezenta
lucrare: același tip de date (StatsBomb), aceeași formulare a problemei
(clasificare multi-clasă în timp real) și aceleași metrici de evaluare
(RPS, Brier Score, ECE).

\textbf{Maddox (2022)} propune un cadru Bayesian generalizabil pentru
predicția win probability, testat pe trei sporturi diferite: fotbal
american (NFL), baschet (NBA) și hochei pe gheață (NHL)
\cite{Maddox2022}. Seturile de date utilizate sunt publice, provenite
din sursele oficiale ale ligilor respective, acoperind mai multe sezoane
complete. Modelul Bayesian propus obține rezultate competitive față de
modelele discriminative clasice în toate cele trei sporturi, cu avantajul
suplimentar că oferă intervale de credibilitate explicite pentru
probabilitățile estimate, nu doar valori punctuale. Autorii concluzionează
că un cadru probabilistic unificat poate captura incertitudinea mai bine
decât modelele discriminative și că abordarea se transferă natural între
sporturi cu structuri de joc diferite. Relevanța pentru prezenta lucrare
constă în argumentul privind importanța calibrării probabilistice față de
simpla acuratețe de clasificare, argument care motivează alegerea metricilor
de evaluare utilizate.

\textbf{Hill (2022)} adaptează modelele de win probability pentru fotbalul
canadian (CFL), un sport cu structură similară fotbalului american, dar cu
reguli diferite \cite{Hill2022}. Autorul utilizează date play-by-play
din șase sezoane CFL (2014--2019), reprezentând aproximativ 500 de meciuri,
construind caracteristici temporale similare celor din fotbalul american:
scor curent, timp rămas, poziția pe teren și posesia balonului. Modelul
final, bazat pe regresie logistică cu caracteristici inginerizate, obține
un RPS de
0.171 și o acuratețe de 76.3\% pe setul de test. Autorul concluzionează
că metodele dezvoltate inițial pentru fotbalul american sunt transferabile
în fotbalul canadian cu ajustări minore și că RPS este metrica cea mai
informativă pentru evaluarea calității probabilistice în acest context.
Relevanța pentru prezenta lucrare este dublă: confirmă transferabilitatea
metodelor de win probability între sporturi și justifică utilizarea RPS
ca metrică principală de evaluare.

\textbf{Burke (2010)} este unul dintre pionierii win probability în timp
real, cu aplicare în fotbalul american (NFL) \cite{Burke2010}. Autorul
utilizează date play-by-play din mai multe sezoane NFL, disponibile public,
și construiește un model bazat pe regresie logistică care estimează
probabilitatea de victorie la fiecare snap al mingii, pe baza scorului
curent, timpului rămas, yard-urilor până la primul down și poziției pe
teren. Modelul obține o acuratețe de aproximativ 80\% și un Brier Score
competitiv față de liniile de bază disponibile la acea dată. Concluzia
centrală a autorului este că fiecare eveniment dintr-un meci poate fi
cuantificat prin impactul său marginal asupra probabilității de victorie,
introducând astfel conceptul operațional de Win Probability Added (WPA).
Deși contextul sportiv diferă față de fotbal, principiul fundamental,
actualizarea continuă a probabilității pe baza evenimentelor, este
identic cu abordarea din prezenta lucrare, iar WPA stă la baza
algoritmului de detectare a momentelor decisive implementat.

\textbf{Tango et al. (2007)} introduc conceptul de \textit{Win Probability
Added} (WPA) în contextul baseball-ului, într-o lucrare de referință
pentru analiza cantitativă a sportului \cite{Tango2007}. Autorii utilizează
date istorice din Major League Baseball acoperind zeci de sezoane, calculând
probabilitatea de victorie pentru fiecare situație posibilă de joc (inning,
scor, număr de outuri și baze ocupate). WPA este definit ca diferența dintre
probabilitatea de victorie înainte și după un eveniment, cuantificând astfel
contribuția fiecărei acțiuni la rezultatul final al meciului. Autorii
concluzionează că WPA oferă o modalitate obiectivă de a identifica momentele
decisive ale unui meci și de a evalua contribuția individuală a jucătorilor
dincolo de statisticile clasice. Relevanța pentru prezenta lucrare este
directă: algoritmul de detectare automată a momentelor decisive este
construit pe același principiu, identificarea variațiilor semnificative
ale win probability de-a lungul meciului, adaptat însă pentru fotbal și
pentru date la nivel de eveniment StatsBomb.